\section{本章小结}
\label{sec:ellip-summary}

开篇的问题是：平衡场的每一点由什么决定？本章的回答是，由整个区域内的源项和整条边界上的数据共同决定，这种全局依赖体现在解的每一种表示之中。

\paragraph{解的结构}
区间上的两次积分展示了最基本的结构：边界贡献和源项贡献线性相加，积分常数由边界条件确定。Green函数把源项的贡献写成点源响应的叠加，$G=\Phi-h$把点源的奇性和边界的修正分开；差分矩阵的逆是它的离散版本。谱展开把算子对角化，$u_k=f_k/\lambda_k$说明解比数据多两阶的光滑性，零特征值则要求先处理核空间。这些公式依赖几何或边界的特殊结构，适用范围有限，也都要求足够的光滑性。

\paragraph{弱理论}
弱形式把一阶导数分给测试函数，让系数跳变和粗糙的源都能进入问题。本质边界条件写进空间，自然边界条件写进积分。Lax--Milgram定理把适定性归结为有界、强制和右端连续三个条件，其证明由Riesz表示和强制性下界完成。三类边界条件的区别，在于怎样处理常数方向，见表~\ref{tab:ellip-summary}。

\begin{BookInlineTable}
  \centering
  \caption{三类边界条件的适定性（$\Omega$有界连通，$\bm{A}$一致椭圆）。}
  \label{tab:ellip-summary}
  \begin{tabular}{lll}
    \toprule
    条件 & 如何排除常数 & 可解性与唯一性\\
    \midrule
    Dirichlet & 零迹与Poincaré不等式 & 任意数据，唯一\\
    Neumann & 零均值子空间 & 需相容条件，差常数\\
    Robin & 边界代价，$\int\beta>0$ & 任意数据，唯一\\
    \bottomrule
  \end{tabular}
\end{BookInlineTable}

\paragraph{三种进一步的结构}
对称问题的弱解是能量的唯一极小点，能量差恰好是误差在能量范数下的平方；非对称问题仍可用Lax--Milgram定理，但没有这种极小刻画。最大值原理依靠符号结构控制逐点的大小，$c\geq0$是必要的结构条件。正则性在内部由系数和源项决定，整体还受边界几何和边界条件切换点的限制。

\paragraph{结论的适用范围}
能量方法和Lax--Milgram定理给出弱解的存在和唯一，对边界和系数的要求很低。最大值原理要求符号条件，经典形式要求解是$C^2$的，弱形式的版本用截断测试函数绕过了这一点。更高的光滑性要求系数、源项和边界同时配合。所有这些结论都是关于平衡场的，没有时间的参与。

第~\ref{chap:diffusion-equations}章在同样的空间和边界模式上引入时间，第~\ref{chap:pde-numerical-methods}章把这里的弱形式在有限维子空间上离散，得到有限元方法。
